\section*{Chapter \ref{ch:b3:sequence-stratigraphy} --- Sequence Stratigraphy}

\begin{solution}{exo:b3:sequence-stratigraphy:1}
$A = 0.5 + 0.1 = \qty{0.6}{m/kyr}$; in \qty{20}{kyr}, \qty{12}{m} of
accommodation. With the sea falling, $A = -0.3 + 0.1 = \qty{-0.2}{m/kyr}$: a
relative fall, slowed by the subsidence.
\end{solution}

\begin{solution}{exo:b3:sequence-stratigraphy:2}
(a) Falling-stage tract (\omterm{def:b3:sequence-stratigraphy:regression}{forced regression}). (b) Transgressive tract. (c)
Highstand tract. (d) Lowstand tract (\omterm{def:b3:sequence-stratigraphy:regression}{normal regression} after the fall).
\end{solution}

\begin{solution}{exo:b3:sequence-stratigraphy:3}
Each \omterm{def:b3:sequence-stratigraphy:parasequence}{parasequence} ends shallower than the one below: a progradational stack,
the shoreline moving seaward, as in a highstand or a lowstand tract. The
opposite order, each ending finer and deeper, is retrogradational: a
transgressive tract.
\end{solution}

\begin{solution}{exo:b3:sequence-stratigraphy:4}
Five couplets to a bundle is the ratio of the short eccentricity (about
\qty{100}{kyr}) to the precession (about \qty{20}{kyr}): each couplet is a
precession cycle, each bundle an eccentricity cycle. The rate is
$1.6/100 = \qty{1.6}{cm/kyr}$, and a couplet of \qty{0.32}{m} lasts about
\qty{20}{kyr}, consistent.
\end{solution}

\begin{solution}{exo:b3:sequence-stratigraphy:5}
$\mathrm dx_s/\mathrm dt = (4000 - 0.15\times20\,000)/50 = \qty{20}{m/kyr}$,
seaward. The shoreline stands still when $AW = Q$: $A = 4000/20\,000 =
\qty{0.2}{m/kyr}$.
\end{solution}

\begin{solution}{exo:b3:sequence-stratigraphy:6}
$\mathrm dx_s/\mathrm dt = 4000/50 + 0.2/0.01 = 80 + 20 = \qty{100}{m/kyr}$;
the fall alone accounts for \qty{20}{m/kyr}, a fifth.
\end{solution}

\begin{solution}{exo:b3:sequence-stratigraphy:7}
It retreats when $Q < AW = 0.4\times25\,000 = \qty{10000}{m^3/kyr}$ per
metre. With $Q = 6000$: $(6000 - 10\,000)/10 = \qty{-400}{m/kyr}$, a retreat
of \qty{0.4}{km} every thousand years: a low shoreface makes a fast
transgression.
\end{solution}

\begin{solution}{exo:b3:sequence-stratigraphy:8}
$1/94.9 - 1/123.9 = \num{0.00247}\text{ cycles}/\unit{kyr}$: a record longer than
$1/0.00247 = \qty{406}{kyr}$ separates them, and \qty{600}{kyr} (resolution
0.00167) is enough. $1/94.9 - 1/98.9 = 0.00042$: it needs about
\qty{2400}{kyr}, four times longer than the record.
\end{solution}

\begin{solution}{exo:b3:sequence-stratigraphy:9}
$16.4/(20\times40.98) = \qty{0.0200}{m/kyr}$, \qty{2.0}{cm/kyr}; the ash beds
are $3.1/0.020 = \qty{155}{kyr}$ apart.
\end{solution}

\begin{solution}{exo:b3:sequence-stratigraphy:10}
(a) $\tan\alpha = -0.4/200 = -0.002$: descending at \ang{0.11}. (b)
$0.2/50 = 0.004$: climbing at \ang{0.23}. A seismic section spans tens of
kilometres at a vertical exaggeration of tens to hundreds, so the angle shows;
a cliff a few hundred metres long shows a rise of a metre or less, lost in
the variation of the beds.
\end{solution}

\begin{solution}{exo:b3:sequence-stratigraphy:11}
$1\,296\,000/54.725282 = \qty{23.68}{kyr}$ ($g_5$); $/57.924383 =
\qty{22.37}{kyr}$ ($g_2$); $/68.377912 = \qty{18.95}{kyr}$ ($g_4$);
$/67.834313 = \qty{19.11}{kyr}$ ($g_3$). $g_2 - g_5 = 3.199101''$ per year,
\qty{405.1}{kyr}; $g_4 - g_5 = 13.65263''$, \qty{94.9}{kyr}.
\end{solution}

\begin{solution}{exo:b3:sequence-stratigraphy:12}
Both make $0.3\times100 = \qty{30}{m}$ of accommodation at the same rate; with
the same supply the geometry is the same, since sediment responds only to the
relative sea level. A eustatic change is the same everywhere: the same
sequence, of the same age, on other margins points to the first history. A
subsidence history built from the basin's thermal and load history
(backstripping), or an independent eustatic record such as oxygen isotopes,
separates them too.
\end{solution}

\begin{solution}{pb:b3:sequence-stratigraphy:1}
\textbf{1.} $\dot E = 30\times2\pi/400\times\sin(2\pi t/400)$: at most
\qty{0.471}{m/kyr}, rising at \qty{100}{kyr} and \qty{500}{kyr}, falling at
\qty{300}{kyr}.

\textbf{2.} \qty{0.020}{m/kyr} at the hinge, 0.037 at \qty{20}{km}, 0.062 at
\qty{50}{km}, \qty{0.120}{m/kyr} at \qty{120}{km}.

\textbf{3.} $A < 0$ when $\sin(2\pi t/400) < -0.0346/0.471 = -0.0734$: $t =
200 + (400/2\pi)\arcsin0.0734 = \qty{204.7}{kyr}$.

\textbf{4.} With $\sigma = 0.0617$: $t = 400 - 63.66\arcsin0.131 =
\qty{391.6}{kyr}$.

\textbf{5.} About \qty{187}{kyr} against 200: the eustatic fall must first
overtake the subsidence, and stops beating it before the lowstand.

\textbf{6.} $0.045\times400 = \qty{18}{m}$, against a eustatic range of
\qty{60}{m}: a third.

\textbf{7.} Landward from 50 to about \qty{120}{kyr}, seaward to about
\qty{440}{kyr}, landward to about \qty{560}{kyr}, seaward again to
\qty{600}{kyr} (the model's turning points are 118.5, 439.5 and
\qty{559}{kyr}; a table every \qty{50}{kyr} only brackets them).

\textbf{8.} $(6000 - 0.30\times15\,000)/40 = \qty{37.5}{m/kyr}$, seaward.

\textbf{9.} $6000/15\,000 = \qty{0.40}{m/kyr}$.

\textbf{10.} $6000/35.7 + 0.424/0.012 = 168.1 + 35.3 = \qty{203}{m/kyr}$,
within \qty{0.5}{\percent} of the model.

\textbf{11.} The supply: \qty{83}{\percent} of the speed.

\textbf{12.} $(-42.7 - 16.3)/(50.7 - 17.4) = \qty{-1.77}{m/km}$, \ang{0.10}.
Since its time, the first point has subsided $0.0345\times395.5 =
\qty{13.6}{m}$, the second $0.0623\times208.5 = \qty{13.0}{m}$: nearly the
same, so in present elevations the descent, \qty{59.0}{m}, differs from the
eustatic fall only by \qty{0.6}{m}. The relative fall at the shoreline at
the time, eustasy minus the subsidence during the fall, was smaller.

\textbf{13.} 150 and 200: highstand; 250, 300, 350: falling stage; 400:
lowstand; 450, 500, 550: transgressive (the shoreline retreats just after
440); 600: highstand.

\textbf{14.} 204.7 and \qty{391.6}{kyr} against 204.5 and 391.5: the same,
within the model's step of \qty{0.5}{kyr}.

\textbf{15.} On the shelf, the surface exposed during the fall: the top of
the highstand deposits from the hinge to \qty{17}{km}, continued over the
downstepping falling-stage foresets; in the basin, the correlative
conformity at the base or top of the falling-stage tract, by convention. A
real shelf would be cut by rivers into \omterm{def:b3:sequence-stratigraphy:condensed}{incised valleys}, weathered into soils,
and partly eroded.

\textbf{16.} $(18.5 - 52.8)/(559 - 439.5) = \qty{-0.29}{km/kyr}$; the
\omterm{def:b3:sequence-stratigraphy:surfaces}{maximum flooding surface} forms at its end.

\textbf{17.} The shoreline sweeps quickly over a nearly flat coastal plain:
the supply is trapped near it, spread over a wide area, and each place
receives little. Basinward of the toe, beyond about \qty{55}{km}, the sea
floor is starved: a \omterm{def:b3:sequence-stratigraphy:condensed}{condensed section} forms there.

\textbf{18.} Resolution $1/30 = \num{0.033}\text{ cycles}/\unit{m}$; Nyquist $1/(2\times
0.01) = \num{50}\text{ cycles}/\unit{m}$.

\textbf{19.} $7.5/2.0 = 3.75$, 1 and $0.81/2.0 = 0.405$; the orbital ratios
are $405/100 = 4.05$ and $41/100 = 0.41$.

\textbf{20.} \qty{2.0}{m} for \qty{100}{kyr}: \qty{2.0}{cm/kyr}; $0.81/41 =
\qty{1.98}{cm/kyr}$; $7.5/405 = \qty{1.85}{cm/kyr}$. The wavelengths the
\omterm{def:b3:sequence-stratigraphy:periodogram}{periodogram} can show near \qty{8}{m} are 10, 7.5 and \qty{6}{m}: the true
\qty{8.1}{m} falls between two of them.

\textbf{21.} $30/0.020 = \qty{1500}{kyr}$.

\textbf{22.} $23.7\times0.020 = \qty{0.474}{m}$. Sampled every
\qty{25}{cm}, the Nyquist frequency is 2 cycles per metre, a wavelength of
\qty{0.5}{m}: the precession, at 2.11 cycles per metre, is beyond it and
folds back to $4 - 2.11 = 1.89$, a false wavelength of \qty{0.53}{m}.

\textbf{23.} A single peak fits some period at some rate whatever its
origin; three peaks whose ratios match the orbital ratios are independent of
the rate, which must then come out the same from each.

\textbf{24.} The shoreline moves from \qty{17.37}{km} at the start of the
fall to \qty{52.84}{km} at the end of the lowstand regression: a maximum
basinward shift of \textbf{\qty[text-series-to-math=true]{35.5}{km}}.
\end{solution}
